% PAPER 4 — exploratory exact-source rank-cell certificate note.
% NOT a prospective transfer improvement. NOT JOURNAL READY.
\documentclass[11pt,a4paper]{article}
\usepackage[margin=0.95in]{geometry}
\usepackage[T1]{fontenc}
\usepackage{amsmath,amssymb,booktabs,graphicx}
\usepackage[hidelinks]{hyperref}
\usepackage{lmodern}
\usepackage{placeins}
\title{Polynomial Rank Cells for Continuous EEG Shrinkage Paths:\\
A Retrospective Source-Domain Certificate Study}
\author{AUTHORS TO BE VERIFIED}
\date{}
\begin{document}\maketitle
\begin{abstract}
\textbf{Objective.} For a source-trained
regularised EEG linear discriminant, determine
which continuous shrinkage intervals can
be certified not to change any held-source
positive--negative score ordering.
\textbf{Methods.} We express each positive--negative
source score difference as a rational function
of the scalar shrinkage parameter. Its numerator
is a polynomial of degree at most the
feature dimension minus one, giving a finite
rank-cell decomposition on feasible intervals.
A Bernstein-basis enclosure uses whole-numerator
cancellation rather than only separate
termwise spectral enclosures. We test the
two source-only certificate procedures
on the same previously exposed
three-channel Weibo2014 motor-imagery
feature folds without changing any
held-target prediction or source selector.
\textbf{Results.} In nine held-source
validation groups under the original
six-dimensional tangent feature geometry,
the proportion of pair orderings certified
unchanged on the full $[0,1]$ path rose
from $94.121\%$ to $94.267\%$
($+0.146$ percentage points). On intervals
between source choices, coverage rose
from $96.274\%$ to $96.338\%$
($+0.064$ percentage points). A separate synthetic sensitivity check across 15 SPD
spectral states in dimensions $p=2$ to $12$
found coverage improvements ranging from zero
at $p=2$ to more than 16 percentage points on
some higher-dimensional synthetic configurations;
it does not constitute new EEG validation.
A stronger 960-interval exact-rational comparison
certified 912 intervals using unshifted Horner,
but \emph{924 each} with midpoint-centered Horner,
Bernstein or their selective hybrid. The 12
Bernstein-only certificates against the weaker
baseline therefore vanished against centered
Horner; the latter was faster in this
machine-dependent microbenchmark.
All 36 remaining intervals contain exact rational
positive and negative witnesses, so 924/960 is the
strict whole-interval sign-coverage ceiling for
this particular fixture family. A corrected exploratory
4,320-interval rational suite certified 4,232 with
Bernstein versus 4,212 with midpoint-centered Horner,
a twenty-interval gain. The first expansion was
repaired after finding duplicated contrast patterns. An incremental exact event-map prototype also returns
the same roots, tie values and cell AUCs as full pair
re-evaluation, while reducing aggregation time on
event-rich synthetic cases; no stable gain appears on
no-event controls. This is not a floating EEG pipeline proof.
\textbf{Conclusion.} Whole-numerator
bounds certify source rank constancy, but
a small Bernstein advantage over
midpoint-centered Horner appears only
in a small subset of the corrected synthetic conditions.
No exclusive state-of-the-art numerical
method, end-to-end floating EEG certificate,
or prospective transfer gain is established.
\end{abstract}
\noindent\textbf{Keywords:} EEG;
regularised discriminant analysis;
spectral perturbations; polynomial
certification; rank stability;
Bernstein basis.

\section{Introduction}
Regularisation is often selected by
cross-validation on source participants.
A linear classifier's held-source
ROC-AUC depends only on positive--negative
pair orderings, so a continuous scalar
shrinkage path need not produce a
continuous performance objective.
The source score is a rational function
of shrinkage under the classical
covariance eigendecomposition. Instead of
optimising its held-target outcomes,
we quantify how much of the
SOURCE AUC path can be certified
stable by algebraic sign enclosures.
Riemannian EEG decoding and source
evaluation follow established methods
\cite{barachant2012,jayaram2018moabb}.
This is distinct from Paper 1's
Gaussian-vs-discriminative measurement
question, Paper 2's recording-content
identity audit, and Paper 3's
held-target policy efficacy question.

\section{Methods}
For fitted source covariance
$S=U\operatorname{diag}(s_i)U^\top$
and source class contrast $d$,
the score direction satisfies
$w_\lambda=[(1-\lambda)S+
\lambda\nu I]^{-1}d$ with
$\nu=\operatorname{tr}(S)/p$.
For any fixed held-source
positive--negative feature difference $v$,
the sign of $v^\top w_\lambda$
changes only where its
rational numerator vanishes or
where the denominator leaves the
strictly positive feasible region.
The following technical formulation
retains this sign convention, treats ties
separately and avoids false formal
machine-verification claims.

\subsection{Finite rank cells on the continuous spectral shrinkage path}
\label{sec:v25-bernstein-rank-cells}
The V23 interval method bounded source-trial score signs by
enclosing the individual rational terms. In low-dimensional
EEG, a complementary approach acts on their \emph{common
numerator}, exploiting cancellations otherwise lost in
termwise bounds.

For fixed source-trained eigensystem
$S=U\,\mathrm{diag}(s_1,\ldots,s_p)U^\top$, class contrast
$d_s$, fixed held-\emph{source} positive--negative feature
difference $v_{pn}$, let
\[
q_i(\lambda)=(1-\lambda)s_i+\lambda\nu,\qquad
a_{i,pn}=(u_i^\top v_{pn})(u_i^\top d_s).
\]
On a feasible interval $I$ where $q_i(\lambda)>0$ for all
$i,\lambda$, the sign of the score difference is the sign
of the single polynomial
\begin{equation}
 N_{pn}(\lambda)
 =\sum_{i=1}^{p}a_{i,pn}\prod_{j\ne i}q_j(\lambda),
 \qquad\deg N_{pn}\le p-1 .
 \label{eq:v25-spectral-numerator}
\end{equation}
Indeed $z_{pn}=N_{pn}/\prod_iq_i$ and the denominator is
strictly positive. Unless $N_{pn}$ vanishes identically
(and the pair remains tied throughout), a positive--negative
pair can change its strict ordering only at the at most
$p-1$ distinct real roots of $N_{pn}$. Hence the finite
source-validation macro-AUC is piecewise constant on
the complement of the finite union of all such rank-cell
boundaries. Values \emph{at} a root are evaluated separately
because ties count for one-half AUC credit; it would be
incorrect to assume the supremum is attained only at
open-cell midpoints. The root-count argument is an
elementary consequence of rational functions, not a
claim of priority for new polynomial algebra.

For each subinterval $[\ell,r]$, map $\lambda$ to
$t\in[0,1]$ and express the degree-$p-1$ numerator
as a Bernstein polynomial
\[
 N_{pn}(\ell+(r-\ell)t)
 =\sum_{k=0}^{p-1} b_{k,pn}B_k^{p-1}(t).
\]
As Bernstein basis functions are nonnegative and sum to
one, $\min_k b_{k,pn}\le N_{pn}(\lambda)\le
\max_k b_{k,pn}$. Strict positivity (or negativity)
of \emph{all} coefficients proves that the pair's rank
is unchanged on the entire subinterval. A source pair
is declared rank-stable over $I$ only if \emph{every}
subinterval proves the same strict sign. A failure to
certify is retained as uncertainty, not counted as a
real crossing. The inference is exact in real arithmetic;
our floating-point implementation applies a conservative
cancellation margin but does not use formally verified
directed rounding, and therefore we do not claim a
machine-checked mathematical certificate.

The new implementation uses a positive product-of-denominators
scaling and is explicitly restricted to $p\le12$ to
control numerical conditioning. With three EEG channels,
a tangent-space covariance model has only six unique
coordinates, yielding a numerator of degree at most
five per positive--negative sample pair. AUC credit of
all certified source pairs remains unchanged over the
entire candidate path. Adding one-sided worst-case
credit changes for the uncertified pairs provides a
deterministic interval-wide finite-\emph{source}-AUC
envelope analogous to the established source macro-AUC bounding construction.

\paragraph{Research scope and falsification.}
The practical hypothesis is that whole-numerator
Bernstein bounds reduce unnecessary uncertainty relative
to the V23 termwise/Taylor interval method on
\emph{identical frozen} held-source model folds and
lambda intervals. We test this on a previously exposed
three-channel Weibo2014 acquisition as a retrospective
mechanistic comparison, not as a fresh target-validation
sample. Either method may be tighter; gains in certified
\emph{source} rank coverage are not gains in unknown-person
AUC, and no claim of cross-subject error prevention
is implied. Boundary roots, exact ties, denser intervals,
numerical cancellation, and subject-specific harms remain
necessary negative controls.


% Exact rational arithmetic audit: computational proof scope, not new EEG evidence.
\subsection{Exact-arithmetic confirmation of the certificate implication}
\label{sec:p4-exact-rational}
To separate a mathematically valid enclosure from the numerical
implementation used for the retrospective EEG experiment, we
implemented an additional independent rational-input verifier.
For rational $s_i>0$, $\nu>0$ and $a_i$, it forms
$N(\lambda)=\sum_i a_i\prod_{j\ne i}
(s_j+(\nu-s_j)\lambda)$ with exact rational coefficients.
On a rational interval $[l,r]\subseteq[0,1]$, it performs
an exact affine substitution followed by the power-to-Bernstein
basis conversion
\[
 b_k=\sum_{j=0}^{k}c_j
 {\binom{k}{j}\over\binom{n}{j}},\qquad
 N(l+(r-l)t)=\sum_{j=0}^{n}c_jt^j.
\]
A strictly positive (negative) set of \emph{all} rational
Bernstein coefficients implies the same strict sign at
\emph{every} point on the entire closed interval because
$B_k^n(t)\ge0$ and $\sum_kB_k^n(t)=1$.
The all-zero polynomial is a separate identically tied case;
mixed signs are \emph{unresolved}, not automatically an actual
rank reversal. This sufficiency statement needs neither a
dense grid nor random interpolation points.

The executable independent check exercises rational spectra in
dimensions $p=2,3,6,9,12$, with cancellation patterns,
boundary roots and exact ties. Its purpose is to test the
logical and arithmetic implementation of this narrow certificate,
NOT to revise real-EEG pair-coverage estimates. In particular,
the original eigenvalues and score projections come from
floating-point preprocessing and are not exact rational
measurements of the underlying physical signal. Exact
calculations on rationalized inputs alone would not certify
that preprocessing, eigendecomposition or source feature
estimation were numerically error-free. A fully verified
real-EEG guarantee still requires input uncertainty bounds
and directed rounding or an independently checked interval
arithmetic implementation.

\subsection{Centered-Horner enclosure and selective Bernstein refinement}
\label{sec:p4-selective-hybrid-proof}
To obtain a meaningful competitor for the original
Bernstein-only enclosure, we add a midpoint-centered,
exact-rational interval Horner evaluator and a selective
two-stage certificate. Let the exact rational score-numerator
polynomial be $P(\lambda)=\sum_{j=0}^d a_j\lambda^j$
on $I=[\ell,r]\subseteq[0,1]$, with
$c=(\ell+r)/2$ and $h=(r-\ell)/2$. The translated
polynomial in $u=\lambda-c\in[-h,h]$ is
\[
P(c+u)=\sum_{k=0}^d b_k u^k,\qquad
b_k=\sum_{j=k}^d \binom{j}{k}a_jc^{j-k}.
\]
An interval Horner evaluation of these $b_k$
over $[-h,h]$ encloses every $P(\lambda)$ on $I$
by interval-arithmetic inclusion. Since the input
coefficients and endpoints are rational, our exact
\texttt{Fraction} arithmetic introduces no numerical
rounding at this step. Neither this fact nor
the power-basis coefficient identity certifies
the unverified floating EEG feature extractor.

\paragraph{Selective refinement algorithm.}
Evaluate centered Horner first. If its lower bound is
strictly positive or its upper bound strictly negative,
retain that sign without computing Bernstein
coefficients. If centered Horner is unresolved,
transform to the Bernstein basis on the \emph{same}
interval and accept a strict sign only if all its
coefficients share that strict sign. Otherwise
return unresolved. A polynomial whose coefficients
all vanish is a separately recorded exact tie.
Inclusion is mandatory at \emph{both}
stages; finite point samples cannot substitute
for either interval proof.

\paragraph{Certificate-union proposition.}
Suppose both the centered-Horner and Bernstein
procedures return a strict sign only when their exact
enclosures exclude zero. Then the selective method
is sound and certifies precisely the union of
the intervals certified by either procedure.
Indeed, a strict centered-Horner sign gives the
desired inclusion proof immediately. Otherwise the
algorithm invokes the Bernstein procedure, which
gives an independent inclusion proof if it succeeds.
If both procedures could certify opposite strict
signs on the same interval, the two sound
enclosures would force $P$ to be simultaneously
strictly positive and strictly negative on a
nonempty interval, which is impossible. No
unknown-sign interval is promoted to a certificate.
Consequently, for a fixed test collection $\mathcal I$,
\[
C_{\rm selective}(\mathcal I)
= C_{\rm centered}(\mathcal I)
\cup C_{\rm Bernstein}(\mathcal I),\qquad
|C_{\rm selective}|\geq
\max\{|C_{\rm centered}|,|C_{\rm Bernstein}|\}.
\]
This is a correctness property, \emph{not}
a theorem of runtime improvement. The hybrid
pays the centered-Horner cost for every interval
and the Bernstein cost on unresolved ones;
when centered Horner rarely resolves or
Bernstein is costly, this policy can be slower
than either stand-alone method.

The method and signed-root/tie tests are archived
in the public experiment source accompanying
workflow \texttt{38326786180}.
The strict mathematical scope is rational-input
polynomial sign determination. Certifying the
floating-point-generated EEG spectrum, tangent
features and original SPD eigensystem would
require further upstream error bounds.

% Exact rational coefficient-only uncertainty; independently executed fixed-seed test.
\subsection{Bounding coefficient uncertainty without assuming a floating EEG certificate}
\label{sec:p4-coeff-uncertain}
For fixed \emph{exact} positive spectrum $s_i>0$, let
$q_i(\lambda)=(1-\lambda)s_i+\lambda\nu$ with
$\nu=p^{-1}\sum_i s_i$. Suppose the pair numerator
coefficients are not exactly $a_i$ but $a_i+\delta_i$ with
\emph{externally justified} independent bounds
$|\delta_i|\leq\varepsilon_i$. Write
\[
B_i(\lambda)=\prod_{j\neq i}q_j(\lambda),\quad
N(\lambda)=\sum_i a_iB_i(\lambda),\quad
E(\lambda)=\sum_i\varepsilon_iB_i(\lambda).
\]
All $q_i$ and $B_i$ are positive for $\lambda\in[0,1]$.
For every admissible perturbation and every parameter,
\[
N(\lambda)-E(\lambda)\ \leq\
\sum_i(a_i+\delta_i)B_i(\lambda)\
\leq N(\lambda)+E(\lambda).
\]
A strict positive Bernstein coefficient vector for
$N-E$ on a specified closed interval proves a robust
positive score difference there; a strict positive
Bernstein vector for $-(N+E)$ proves a robust negative
difference. Otherwise the algorithm returns
\texttt{UNKNOWN}, and must not assert an ordering.
This follows from the convex-hull property of the
Bernstein representation, without relying on sampling.
The rule is sufficient, not complete: \texttt{UNKNOWN}
may contain either a truly stable or an unstable sign.

A self-contained exact-rational implementation checks invalid
spectra, negative error bounds and accidental binary-float
inputs fail closed. In a fixed-seed set of 96 \emph{synthetic}
dimension-2--5 instances it classified 31 positive,
39 negative and 26 unknown, and tested 3,120
independent box-corner/parameter combinations without
contradicting a classified sign. This sampling is a
regression check, not a proof of coverage or tightness.
The uncertainty model covers only pair-score coefficients,
with fixed exact spectral poles. No bound is yet established
for EEG filtering, input digitisation, covariance estimation,
eigenvalues/eigenvectors, or IEEE-754 rounding. Accordingly
the experiment does not license an end-to-end
machine-certified EEG provenance or transfer claim.

\subsection{Joint spectral- and coefficient-box sign certification}
\label{sec:p4-joint-box}
The previous coefficient-only certificate assumes exact spectral poles.
We additionally consider a strictly positive \emph{input box}
$s_i\in[\underline s_i,\overline s_i]$ and an independent
coefficient box $a_i\in[\bar a_i-\varepsilon_i,\bar a_i+\varepsilon_i]$.
Define $\nu=p^{-1}\sum_i s_i$, and for
$\lambda\in[L,R]\subseteq[0,1]$ let
$q_i(\lambda)=(1-\lambda)s_i+\lambda\nu$.
Because all spectral coordinates are positive, every $q_i$ is positive.

At each parameter endpoint $t\in\{L,R\}$ the bounds
\[
(1-t)\underline s_i+t\underline\nu
\leq q_i(t)\leq
(1-t)\overline s_i+t\overline\nu,\quad
\underline\nu=p^{-1}\sum_j\underline s_j,\quad
\overline\nu=p^{-1}\sum_j\overline s_j
\]
hold despite correlations among the $s_i$.
The affine function $q_i(L+(R-L)u)$ has a degree-one Bernstein
representation whose two controls lie in these bounded intervals.
Multiplying these interval Bernstein polynomials with the
positive combinatorial weights
\[
{m\choose i}{n\choose j}/{m+n\choose i+j}
\]
and accumulating interval-valued coefficients yields sound
enclosures $[\underline b_k,\overline b_k]$ for the degree-$(p-1)$
numerator $N(u)=\sum_i a_i\prod_{j\neq i}q_j(u)$.
If every $\underline b_k>0$, then for \emph{all} admissible
spectra, coefficients and parameters, the underlying rational
pair score difference is positive. If every
$\overline b_k<0$, it is negative. Otherwise return
\texttt{UNKNOWN}. This conservative test does not assume the
eigenvalues are independent in the true covariance; it encloses
a larger independent box and may refuse valid instances.

An implementation based solely on Python exact rational
arithmetic completed 126 seeded synthetic instances with
$p=2$--$5$: 52 positive, 50 negative and 24 unknown
certifications. A separate exact-rational point oracle evaluated
208,560 admissible spectral/coefficient box corners and parameter
probes without a contradictory classification. The oracle is
only a regression test; soundness follows from interval
inclusion and the Bernstein convex-hull property. A negative
or positive box result is valid only if the \emph{input}
bounds are valid. Neither this check nor its fixed-pole
predecessor derives intervals from sample noise, numerical
eigendecomposition, eigenvector motion or IEEE-754 rounding.
No end-to-end EEG safety, external predictive improvement,
or novel root-isolation theorem follows from these experiments.

\subsection{Adaptive certified sign or rational falsification on spectral boxes}
\label{sec:p4-rational-adaptive-counterexamples}
The preceding joint spectral--coefficient box is intentionally
overinclusive. We now consider two sound procedures that operate
on exactly the same \emph{input} boxes: (i) a whole-interval
Bernstein enclosure of the rational numerator and (ii) a
whole-interval termwise rational bound, each using exact fractions.
If either returns a strict sign, no sampling is needed to
justify its claim. When both are indeterminate, the refinement
procedure bisects the continuous regularisation interval at
rational midpoints and requires \emph{all} child boxes to certify
the \emph{same} strict sign. It may also bisect the widest
normalised uncertain eigenvalue interval, again requiring
all descendants to agree; the original box is the exact
union of the two sub-boxes. Therefore an accepted parent sign
is sound by induction on the refinement tree.

If refinement remains undecided, an independent rational
point oracle searches admissible corners of the eigenvalue
and coefficient boxes at a fixed grid of rational values of
$\lambda$. Opposite exact signs at two admissible points,
or a single admissible exact zero, are \emph{witnesses that a
strict invariant sign is impossible}. Mere absence of an
opposite-sign sample is \emph{not} a proof of robustness,
so the algorithm otherwise reports \texttt{UNKNOWN}.

In a fixed-seed \emph{constructed} suite with dimensions
$p=2$--$5$, comprising 126 earlier spectral-box fixtures
and 96 additional mixed-sign fixtures (all retained),
whole-interval Bernstein proved 173/222 cases, whole-interval
direct rational bounds 174/222, and $\lambda$ refinement
184/222. Additional independent spectral splitting proved
two further cases. The resulting 186 cases consist of
95 robust positive and 91 robust negative signs.
For the other 36, actual valid exact rational zero/opposite-sign
witnesses refute strict invariance. Consequently the fixture
suite has zero unclassified cases, but this is
\emph{not} a complete algorithm for arbitrary coefficient
boxes and not an independent EEG result.
The machine-readable record stores all 36 falsifying
examples, with exact numerator/denominator values, and
the scripts regenerate each fixture and check admissibility
and sign with a separate rational score formula.
The all-same-sign corner regression for the 186 certified
fixtures is supportive only; formal soundness derives from
positivity of rational denominators, interval enclosure and
the Bernstein convex-hull property.

This exercise illustrates an important negative control:
a greater numerical sign-coverage ratio would be misleading
if exact admissible counterexamples already preclude the
stability property. We do not attribute a new optimality theorem,
IEEE-754 rounding control, eigenvector motion bound or
improvement of EEG transfer AUC to this constructed result.

\subsection{Covariance-coordinate sign robustness under eigenspace rotation}
\label{sec:p4-cov-resolvent}
The spectral box method above requires \emph{externally justified}
bounds on coefficients expressed in a fixed eigenspace.
A complementary sufficient condition works directly in original
covariance coordinates, including arbitrary rotations of eigenvectors.
Let $S=S^\top\succ0$, $d$ be the fitted class contrast,
$v$ a held-source positive-minus-negative feature difference,
and $\nu=\operatorname{tr}(S)/p$. For any
$\lambda\in[L,R]\subseteq[0,1]$, write
$A_\lambda=(1-\lambda)S+\lambda\nu I$ and
$f_\lambda=v^\top A_\lambda^{-1}d$.
Consider arbitrary symmetric covariance perturbations
$\|\Delta S\|_2\le\epsilon_S$ and feature perturbations
$\|\Delta v\|_2\le\epsilon_v$,
$\|\Delta d\|_2\le\epsilon_d$, all jointly admissible,
and set $\nu'=\operatorname{tr}(S+\Delta S)/p$.
Because
\[
\|A'_\lambda-A_\lambda\|_2
\le(1-\lambda)\|\Delta S\|_2+
\lambda|\operatorname{tr}(\Delta S)|/p
\le\epsilon_S,
\]
the perturbation of the isotropic shrinkage target
is included automatically.

An exact lower bound $\mu_S\le\lambda_{\min}(S)$
is available from the rational inverse row-sum norm:
$\mu_S=1/\|S^{-1}\|_\infty$;
symmetry ensures $\|S^{-1}\|_2\le\|S^{-1}\|_\infty$.
Set $\mu=\min(\mu_S,\nu)$ and assume
$0\le\epsilon_S<\mu$.
Then all feasible shrunk matrices remain positive definite,
and the resolvent identity gives the uniform guarantee
\[
|f'_\lambda-f_\lambda|\le
\frac{\epsilon_S\|v\|_2\|d\|_2}
     {\mu(\mu-\epsilon_S)}
+\frac{\|v\|_2\epsilon_d+
       \epsilon_v\|d\|_2+\epsilon_v\epsilon_d}
      {\mu-\epsilon_S}
\equiv B.
\]
The code replaces both vector 2-norms by their
computable rational 1-norm upper bounds, never by
favourable approximations.

For rational input $S,v,d$, the determinant identity
\[
f_\lambda =
\frac{v^\top\operatorname{adj}(A_\lambda)d}
     {\det A_\lambda} =
\frac{N(\lambda)}{D(\lambda)},\qquad D(\lambda)>0
\]
allows rational polynomial Bernstein controls on $[L,R]$.
If every numerator Bernstein coefficient is positive,
$\inf f_\lambda\ge
\min_k b^N_k/\max_k b^D_k$;
the analogous negative bound uses
$-\max_k b^N_k/\max_k b^D_k$.
A strict score margin greater than $B$ proves the
sign for \emph{all} admissible covariance, contrast and
pair-feature perturbations and all parameters in $[L,R]$.
Otherwise the procedure returns \texttt{UNKNOWN}.
This argument does not assume eigenvector alignment,
simple eigenvalues or a fixed eigenbasis; it is a
conservative sufficient condition, not a sharp theorem
about unknown EEG acquisition noise.

In a seeded constructed suite of 120 exact-rational
symmetric-positive-definite matrices of dimensions
2--4, 46 positive and 65 negative cases were certified;
9 returned unknown. Independently reconstructed
2,220 admissible rational covariance/vector perturbation
and parameter probes contradicted none of these signs.
These point probes check implementation rather than
prove the spectral-norm theorem. Crucially, no validated
mapping from digitised EEG, covariance-estimation error,
or floating-point operations to the externally assumed
$(\epsilon_S,\epsilon_d,\epsilon_v)$ is presently available.
Neither the synthetic outcomes nor this classical
resolvent bound proves a new EEG classification
improvement or unqualified state-of-the-art originality.

\subsection{From certified final-feature input errors to covariance-score signs}
\label{sec:p4-final-feature-chain}
The covariance-coordinate theorem previously assumed
an externally supplied matrix-level error radius.
We now derive that radius \emph{deterministically}
from a more elementary, equally conditional assumption:
the \emph{final stored feature matrix} $\widehat X\in
\mathbb Q^{n\times p}$ is related to the hypothetical
unrounded feature matrix $X$ by entrywise
$|X_{ij}-\widehat X_{ij}|\leq\eta$.
Here $\eta$ is an \emph{externally justified bound
on the final features}, \emph{not} a property we
infer from an EDF's raw 16-bit conversion.
No non-linear EEG preprocessing, tangent transform
or numerical matrix decomposition is covered by
this assumption without an additional verified
forward-error analysis.

For two classes, let $H$ be the symmetric
block-diagonal within-class centring projection,
$R=HX$, $\widehat R=H\widehat X$ and
$\rho I$ a fixed identical positive ridge.
Then the pooled within-class covariances are
\[
S=\frac{R^\top R}{n-2}+\rho I,\qquad
\widehat S=\frac{\widehat R^\top\widehat R}{n-2}+\rho I.
\]
Writing $E=X-\widehat X$ gives
$\|HE\|_2\leq\|E\|_F\leq\eta\sqrt{np}$.
By expanding the Gram product,
\[
\|S-\widehat S\|_2\leq
\frac{2\|\widehat R\|_F\eta\sqrt{np}
       +np\eta^2}{n-2}
\equiv\epsilon_S.
\]
For the difference $d$ of two empirical class
means and any positive-minus-negative example
difference $v$, the same entrywise assumption
implies
\[
\|d-\widehat d\|_2\leq 2\eta\sqrt{p}
\equiv\epsilon_d,\qquad
\|v-\widehat v\|_2\leq 2\eta\sqrt{p}
\equiv\epsilon_v.
\]
Replacing all square roots with \emph{upward exact-rational}
approximations gives a machine-checkable rational
outer bound, which is submitted to the
covariance-coordinate resolvent criterion in
Section~\ref{sec:p4-cov-resolvent}.
If it returns a strict sign, that sign holds for
\emph{every} exact feature matrix in the specified
entrywise input box and every $\lambda\in[0,1]$.
No independence assumptions on feature errors,
no fixed eigenvectors and no sampled noise law
are required for this conditional conclusion.

The included synthetic exact-rational test uses
64 seeded two-class, 24-record final-feature matrices
with a common positive ridge. Of these deliberately
well-separated fixtures, 48 have proved positive
score signs and 16 proved negative signs over the
entire continuous $\lambda$ interval and every
admissible entrywise feature perturbation; no
candidate was inconclusive in this \emph{easy}
synthetic subset. Independently reconstructed
admissible feature perturbations produced 576 exact
rational score checks without a contradiction.
A separate, deliberately adversarial \emph{exact
counterexample} perturbs one positive and one
negative record, each by $\eta=0.1$: the nominal
score is positive and an admissible score is
negative across tested shrinkage values; the
algorithm correctly returns \texttt{UNKNOWN}.
This demonstrates that the apparatus can
refuse certification when input uncertainty
permits a genuine decision reversal.
The program recomputes all class means,
within-class residuals, covariance and pair
differences rather than merely injecting
a matrix-level error. The successful code,
exact samples and fail-closed counterexample
are archived in
\url{https://github.com/lindicaphxag-tech/kaggle/actions/runs/38389194383}.
These are conditional operator regressions,
not real EEG classification, a patient-level
risk model, or proof that any unverified
preprocessing/tangent pipeline satisfies
the assumed final-feature $\eta$ bound.

\subsection{Real publicly released EEG: exact rounding bounds on final tangent features}
\label{sec:p4-real-feature-quantisation}
To test the conditional input-to-covariance certification
on \emph{non-synthetic EEG features}, we executed a
new numerical-representation experiment using
\emph{previously analysed} BNCI2014\_001
public released person 1.
Two EEG channels (C3/C4) were band-pass filtered
8--30 Hz and resampled at 160 Hz using
MOABB 1.7.2. The source-only tangent geometry
yielded three real \emph{binary64} final feature
coordinates. A fixed 64-trial (32 per class)
feature matrix and positive ridge $1/5$
were used in an ordinary within-class
covariance estimate. Twelve positive--negative
sample pairs were fixed by a seed
\texttt{20261013}, before any certification
results were inspected; the experiment is
a rank-sign \emph{numerical accuracy}
demonstration and does \emph{not} evaluate
cross-person decoding benefit.

Each final binary64 feature was converted
\emph{exactly} to its IEEE finite-binary
rational value, then rounded to a fixed
rational grid with resolution $2^{-18}$.
The exact original rational and quantized
rational values were compared
\emph{elementwise}, confirming the bound
$\eta=2^{-19}$ without relying on
machine-rounding approximations.
The deterministic covariance- and
contrast-error radii derived in
Section~\ref{sec:p4-final-feature-chain}
were then submitted to the eigenvector-
independent resolvent certificate in
Section~\ref{sec:p4-cov-resolvent}.
For the preselected twelve held pair
comparisons, eight strict positive
and four strict negative score signs
were proved simultaneously for
\emph{all} $\lambda\in[0,1]$ and
\emph{every} rational or real
final-feature matrix inside the
verified $\eta$ input box; none
returned unknown. The complete
public EEG execution, frozen
sample selectors and rational
certificate parameters are archived at
\url{https://github.com/lindicaphxag-tech/kaggle/actions/runs/38391746843}.

This is an actual EEG \emph{final-feature
fixed-point representation test},
but it does not certify that a
true physical EEG acquisition
satisfies any analogous input
bound. The original MNE filtering,
covariance, manifold tangent mapping
and eigenvalue computations
are \emph{outside} the verified
input error box. Nor do 12
sample-pair signs imply a general
classification accuracy, population
transfer advantage, novel root
isolation algorithm, or meaningful
worst-person BCI benefit.
The all-12 result is descriptive
for this deliberately small
three-coordinate previously exposed
recording only.


\section{Retrospective source-only study}
\subsection{Real three-channel EEG source rank cells: a small numerical improvement}
\label{sec:v25-real-bernstein}
An executable retrospective source-only audit was completed on
Weibo2014 using the same original C3/Cz/C4 sensorimotor montage,
fixed six-dimensional tangent-LDA pipeline, and first outer
held-subject fold as the V24 first-run study. The other nine
participants each entered the original held-source validation
procedure, with source-fitted tangent references and shrinkage
spectra unchanged. Because Weibo2014 target outcomes had already
been exposed in V24, this audit is a \emph{post-outcome
mechanistic analysis}; it is not a new, independently untouched
external test or a prospective predictor of target harm.

\begin{table}[tb]
\centering\small
\caption{V25 3-channel Weibo2014 source-domain certificate
comparison, nine held-\emph{source} people, 12 fixed lambda
subintervals per range. Percentages count only pair rankings
conservatively proved invariant along an entire continuous
shrinkage interval. Failure to certify is \emph{not} an
observed rank crossing. There are no new target AUC effects
in this table.}
\label{tab:v25-source-bernstein}
\begin{tabular}{lrrr}
\toprule
Source interval & V23 Taylor & V25 Bernstein & Gain (pp)\\
\midrule
Complete $[0,1]$ & 94.121\% & 94.267\% & +0.146\\
Between source selectors & 96.274\% & 96.338\% & +0.064\\
\bottomrule
\end{tabular}
\end{table}

V25 gave a slightly higher certifiable pair fraction for
all nine validation source people on both prescribed
intervals, and reproduced the same endpoint AUC values
as V23 to the frozen numerical tolerance. These
increases are \emph{small}; no genuine future-person
generalisation effect, clinical value, computational
advantage or superior transfer selector has been
established. Neither method guarantees formal verified
floating-point arithmetic. The original completed
GitHub Actions proof-of-execution is
\texttt{38283087191}, with frozen JSON SHA-256
\texttt{c0c9e038376c899d}\ldots.
The full digest and source-person results remain in
the public immutable experiment artifact
\texttt{11762626131}.

The retrospective result motivates further mathematical
study of spectral rank boundaries, but does not justify
promoting the whole-numerator certificate into the
primary manuscript's EEG decoding performance claims.


\section{Synthetic multi-dimensional numerical experiment}
% Additional independent, fully synthetic stress profile; NOT a new EEG dataset.
\subsection{Cross-dimension numerical stress test, not neural-transfer evidence}
\label{sec:p4-synthetic-multidimensional}
A second, separately executed code path tested the two
source-domain bounding implementations on constructed
symmetric-positive-definite eigensystems of dimensions
$p\in\{2,3,6,9,12\}$. Each setting used three fixed
random seeds, 18 synthetic positive and 18 synthetic
negative feature rows (324 score-pairs per setting),
12 shrinkage subintervals in $[0.02,0.98]$ and 201
directly recomputed $\lambda$ values. In all fifteen
cases, both calculated source AUC envelopes contained
all dense-grid AUC values. This is a sampled numerical
consistency check, \emph{not} a mathematical proof of
floating-point enclosure between sampled grid values.

\begin{table}[tb]
\centering\small
\caption{Synthetic SPD spectral certificate sensitivity
across feature dimension. Mean percentage-point
difference in SOURCE pair ranking stability coverage,
Bernstein minus existing termwise/Taylor, over three
fixed seeds (not independent EEG participants).}
\label{tab:p4-synthetic-spd}
\begin{tabular}{rrrl}
\toprule
Feature $p$ & Trials & Mean coverage gain (pp) &
Bernstein versus old\\
\midrule
2 & 3 & $+0.000$ & 3 ties\\
3 & 3 & $+4.321$ & 3 better\\
6 & 3 & $+13.169$ & 3 better\\
9 & 3 & $+16.667$ & 3 better\\
12 & 3 & $+16.152$ & 3 better\\
\bottomrule
\end{tabular}
\end{table}
These are deliberately generated synthetic spectra,
not additional EEG acquisitions; their coverage
percentages cannot be pooled with or substituted for
the observed Weibo2014 $+0.146$-pp real SOURCE improvement.
The synthetic experiment neither demonstrates a
computational speedup nor identifies how often
uncertain source pairs actually reverse rank. Although
the higher-dimensional means suggest whole-numerator
cancellation may be increasingly useful for these
particular constructed states, synthetic generator
choice and floating-point instability can change
the result. Validation on multiple NEW native
EEG feature geometries remains necessary.
The source code, raw per-case results and all
fifteen dense-grid checks are linked in the
supplementary execution receipt (run 38312807723).

\subsection{Exact-rational benchmark against classical interval Horner evaluation}
\label{sec:p4-exact-two-compare}
We directly compared an exact rational Bernstein
coefficient enclosure against \emph{standard interval
Horner evaluation} using identical rational spectra,
weights and intervals. The competitor is not a
naive point-sampling heuristic; both enclosures are
mathematically sound whenever they certify a strict
sign. Each of $p=2,3,6,9,12$ included
12 fixed deterministic spectral coefficient patterns
evaluated on the same 16 uniform subintervals of
$[0,1]$, giving 960 subinterval tests.
The exact rational calculations do not suffer
from binary floating rounding in the test itself,
but this does \emph{not} certify EEG acquisition or
floating-point eigendecomposition.

\begin{table}[tb]\centering\small
\caption{Actual executed exact-rational comparisons.
Counts are strict-sign certificates out of 192
tested intervals at each dimension. Timings in
seconds are Python/GitHub-runner inner sign-check
measurements only, not measured total pipeline
throughput or proven complexity.}
\label{tab:p4-horner-comparison}
\setlength{\tabcolsep}{4pt}
\begin{tabular}{rrrrrr}
\toprule
$p$ & Bernstein & Horner & Extra & Bernstein s & Horner s\\
\midrule
2  & 180 & 180 & 0 & 0.0087 & 0.0025\\
3  & 192 & 192 & 0 & 0.0120 & 0.0038\\
6  & 180 & 180 & 0 & 0.0409 & 0.0098\\
9  & 192 & 192 & 0 & 0.0988 & 0.0181\\
12 & 180 & 168 & 12 & 0.1630 & 0.0230\\
\midrule
All & 924 & 912 & 12 & 0.3235 & 0.0572\\
\bottomrule
\end{tabular}
\end{table}

Bernstein was able to certify 12 additional
strict-sign intervals, all at dimension 12,
while Horner alone certified none in this
specific deterministic fixture family.
This gain was $12/960=1.25$ percentage points
in the chosen microbenchmark. However, Bernstein
was also slower by about $5.65$ times in its
inner exact-rational sign-check timing
($0.3235$ s versus $0.0572$ s in the
measured run), and the extra gain did not occur
at dimensions 2, 3, 6 or 9.
Thus the data \emph{do not} support a claim
that Bernstein is uniformly more computationally
efficient or universally stronger than
classical interval enclosures. The fixed
benchmark is exploratory and synthetic;
it is not an empirical population of
random real-world numerical algorithms.

The full 960-interval source-level reproducibility
record was generated by the public workflow
\url{https://github.com/lindicaphxag-tech/kaggle/actions/runs/38321261591}.
No original held-target EEG AUC or
source policy was improved by this comparison.

\paragraph{Stronger exact centered-Horner control.}
A follow-up reran the identical 960 rational-input intervals
with an exact Taylor shift to each interval midpoint followed
by exact interval Horner evaluation. Centered Horner and
Bernstein both certified 924/960 intervals, with zero
Bernstein-only and zero centered-Horner-only intervals in this
fixture family. Ordinary interval Horner certified 912/960.
The measured inner-check time for centered Horner was lower
than for Bernstein on this runner; timing is descriptive and
implementation-dependent. The stronger control removes the
apparent coverage advantage in this test and reinforces the
publication HOLD. This is a synthetic exact-input result,
not a certified floating-point EEG pipeline or a proof of
universal method equivalence. The three-way execution record
is \url{https://github.com/lindicaphxag-tech/kaggle/actions/runs/38337557905}.

% Verified real executed rational synthetic benchmark: 38326786180.
\subsection{Stronger centered-Horner baseline falsifies an exclusive Bernstein gain}
\label{sec:p4-strong-centered-control}
The initial 960-interval exact-rational benchmark
found 924 strict-sign certifications for Bernstein
and only 912 for \emph{unshifted} power-basis
interval Horner. We tested whether those
twelve additional certificates survive the
stronger classical baseline of
midpoint-translated interval Horner
(Section~\ref{sec:p4-selective-hybrid-proof}).
All methods used the \emph{same}
$p=2,3,6,9,12$ rational spectra,
12 deterministic patterns per dimension
and 16 fixed subintervals per pattern.
Root-touching, identically tied and
negative-case synthetic tests were
included; an outcome was accepted only
when an exact inclusion enclosure
established a strict sign.

\begin{table}[tb]
\centering\small
\caption{Second, stronger exact-rational matched
960-interval benchmark. Parenthesised times
are in-run Python internal evaluation totals;
they are machine-dependent and not asymptotic
runtime guarantees. The test contains no
original real EEG floating input certification.}
\label{tab:p4-centered-horner}
\setlength{\tabcolsep}{4pt}
\begin{tabular}{lrr}
\toprule
Exact rational interval method & Certified / 960 & Seconds\\
\midrule
Original interval Horner & 912 & 0.0742\\
\textbf{Midpoint-centered interval Horner} & \textbf{924} & 0.2623\\
Bernstein enclosure & 924 & 0.4057\\
Centered Horner then Bernstein & 924 & 0.2753\\
\bottomrule
\end{tabular}
\end{table}

The midpoint-centered evaluator recovered all
12 signs missed by the original Horner baseline,
while Bernstein certified \emph{no additional}
interval beyond centered Horner on these
fixtures. The selective method triggered
Bernstein fallback on only 36/960 intervals,
but none of those fallbacks produced an
additional strict-sign certificate.
Its measured inner-run time was slightly
greater than centered Horner alone
($0.2753$ versus $0.2623$~s),
although less than calling Bernstein on
all 960 intervals. The centered method
was also slower than the weakest,
unshifted original Horner baseline.

These results \emph{falsify} a claim of
exclusive Bernstein superiority over an
appropriately transformed classical
interval enclosure on this dataset.
They do not disprove all possible
Bernstein advantages on other polynomial
families, intervals or computing budgets.
A further method contribution requires
benchmarking against
midpoint-translated Horner,
root isolation, arithmetic-error-bounded
interval methods and numerically stable
coefficient transformations.
The exact executed output and source
version are archived at
\url{https://github.com/lindicaphxag-tech/kaggle/actions/runs/38326786180}.
No real EEG classifier or unseen-person AUC
was improved by this synthetic experiment.

\FloatBarrier
\subsection{Corrected multi-family exact-rational stress test}
\label{sec:p4-multifamily}
We further varied the spectral fixture on a fixed interval grid: three deterministic strictly
positive spectral families (linear, quadratic, and alternating),
five dimensions $p=2,3,6,9,12$, the same twelve distinct-up-to-scale rational
contrast patterns per spectrum, and both eight and sixteen equal
rational subintervals. This produces 4,320 matched interval cases. An initial
contrast generator made its twelve nominal patterns scalar multiples;
we detected this defect and replaced it before this corrected run.
This correction was made after examining the first sweep, so this
expanded comparison is exploratory rather than an independently
preregistered test.
The numerator construction, coefficients, interval endpoints and
both methods use exact rational arithmetic. These are synthetic
classifier spectra, not an end-to-end certificate of the numerical
EEG preprocessing pipeline.

\begin{table}[ht]\centering\small
\caption{Exact rational whole-interval strict-sign certifications
on the corrected 4,320-case exploratory sweep. Opposite-sign rational
point witnesses invalidate a whole-interval constant strict sign for
88 unresolved cases. The extra Bernstein signs include thirteen
quadratic-spectrum dimension-twelve cases; other gains are spread
across the remaining tested spectra and dimensions.}
\label{tab:p4-multifamily}
\begin{tabular}{lr}\toprule
Method or diagnostic & Interval count\\\midrule
Midpoint-centered Horner certified & 4,212\\
Bernstein certified & 4,232\\
Bernstein only versus centered Horner & 20\\
Centered Horner only versus Bernstein & 0\\
Both unresolved with opposite-sign rational witnesses & 88\\
Both unresolved without such a witness & 0\\
\bottomrule\end{tabular}
\end{table}

The Bernstein fallback raises the
strict-sign count from 4,212 to 4,232 on this corrected
synthetic suite, a 0.463-percentage-point change in
coverage (20/4,320). It does not repair the 88 genuinely
sign-changing intervals and does not prove an asymptotic
speed advantage. In the earlier quadratic-spectrum 960-case
fixture, the advantage remains zero. This result narrows the
claim to a small reproducible advantage on selected synthetic
intervals under these corrected patterns; tests of rounding-safe floating implementations,
root isolation, runtime and more diverse polynomials remain
necessary before a numerical-algorithm journal submission.

\subsection{Exact event-map update after root isolation}
\label{sec:p4-incremental-events}
A separate prototype now constructs the complete rational-input
AUC event map. It uses SymPy's existing exact real-root routine
for every pair numerator; the new computational step groups equal
algebraic events and updates the two-class pair credit only for
pairs participating at that event. Odd-multiplicity roots reverse
the strict-order credit, even-multiplicity roots retain it, and
each vanishing pair contributes exactly one-half at the event.
Globally tied pairs and roots at the parameter endpoints remain
separate cases. This is an event-aggregation procedure, not a
new polynomial root isolator.

The exact output (all event points, event AUC values, and open-cell
AUC values) agrees with the original implementation that recomputes
every pair's sign in every cell. The local tests include coincident
events, double roots, full-path ties, and endpoint equality.
A deterministic rational SPD synthetic benchmark used dimensions
$p=3,6$, pair counts 4--32, and 8--64 interior events. At $p=6$
we used a full degree-five numerator; the two in-range roots per
pair were constructed, with the remaining roots outside $[0,1]$.

\begin{table}[ht]\centering\small
\caption{Exploratory local Python timings for identical exact event maps.
Each value is the median of three executions, in seconds. Both
algorithms call the same SymPy exact root routine. The last two
rows have no in-range root and show the event-sweep overhead.}
\label{tab:p4-incremental-events}
\begin{tabular}{rrrrr}\toprule
$p$ & Pairs & Events & Full re-evaluation & Incremental\\\midrule
3 & 16 & 32 & 0.111 & 0.019\\
3 & 32 & 64 & 0.450 & 0.037\\
6 & 8 & 16 & 0.107 & 0.026\\
6 & 16 & 32 & 0.380 & 0.056\\
3 & 32 & 0 & 0.461 & 0.482\\
6 & 16 & 0 & 3.523 & 3.547\\
\bottomrule\end{tabular}
\end{table}

The measured speedup concerns aggregation of an event-rich, exactly
solvable constructed family, and does not establish superiority to
Sturm/Descartes root-isolation implementations, floating-point
interval methods, or realistic EEG polynomial distributions.
For no-event cases the implementation is slightly slower, while
exact root isolation itself dominates. Real-source pair scale,
shared-pole construction costs, precision/rounding envelopes and
memory use remain untested. Executable parity and full timing
receipts are available in the public workflow
\url{https://github.com/lindicaphxag-tech/kaggle/actions/runs/38354399661}.

% Independent CI run 38363154062; exact rational polynomial root-stage scope.
\paragraph{Strong exact root-isolation comparator.}
The incremental event update previously held the exact SymPy
\texttt{real\_roots} engine fixed, comparing only full per-cell
re-evaluation with incident-pair updates. We separately compared that
root engine with SymPy's established exact rational isolating
intervals at width at most $10^{-8}$, using identical constructed
rational pair polynomials in 3 and 6 spectral dimensions. All
intervals contained the corresponding exact algebraic roots with
their multiplicities; an exact root-count check also covered the
interior $\lambda\in(0,1)$. The intervals algorithm is a strong
\emph{existing} comparator, not our contribution.

\begin{table}[h!]\centering\small
\caption{Exact root-stage timing on the public CI machine, median
seconds over four alternating-order repeats. Each method isolates all
real roots of every polynomial, not an end-to-end AUC path.
Times depend on the machine and fixed constructed examples.}
\label{tab:p4-root-isolator}
\begin{tabular}{llrrr}
\toprule
Dimension & Pattern & Pairs & Rational intervals & \texttt{real\_roots}\\
\midrule
3 & Crossing & 16 & .00416 & .01081\\
6 & Crossing & 16 & .00888 & .03562\\
6 & No event & 16 & .04569 & .02819\\
\bottomrule
\end{tabular}
\end{table}

The rational isolator is already faster at the \emph{root stage} in
the crossing examples, while it is slower in the six-dimensional
no-event control. Thus the previous speedup over full AUC
re-evaluation cannot be interpreted as superiority to a strong
root-isolation method. A complete strict algorithm must combine
root isolation, equality and ordering of roots across pairs, tie
scores, incident updates and floating input error control in one
fair end-to-end benchmark. The comparator code and full eight-case
receipt ran in
\href{https://github.com/lindicaphxag-tech/kaggle/actions/runs/38363154062}{the successful independent CI workflow}.

\paragraph{Complete event map with the same strong isolator.}
We next used these exact rational isolating intervals in \emph{both}
complete event-map algorithms. A rational interval that lay wholly
outside $(0,1)$ was rejected without computing its algebraic root;
for any nonrational interval overlapping the feasible parameter range,
the solver falls back to \texttt{real\_roots}. We compared full
per-event pair re-evaluation against incremental updates of incident
pairs. The two outputs agreed exactly with one another and with the
original exact event-map implementation, including endpoint ties,
zero polynomials, coincident roots, even multiplicities and an
irrational-root check.

On the independent CI machine, the six-dimensional 16-pair,
32-event constructed example took .14495~s for full re-evaluation
and .02065~s for the incremental method (approximately 7.0 times
faster), using the \emph{same} stronger root engine. The
three-dimensional 32-pair, 64-event case took .08380 versus
.01856~s (approximately 4.5 times faster). In the six-dimensional
16-pair \emph{no-event} control the times were .04998 and
.05000~s, so there is no meaningful gain. This isolates a
conditional event-update advantage on these constructed rational
inputs; it does not imply faster roots, a universal speedup,
floating-input correctness, or a real-EEG performance effect. The
complete six-case workflow is
\href{https://github.com/lindicaphxag-tech/kaggle/actions/runs/38364906164}{available with an exact parity check}.


% P4 executed GitHub Action 38315929126; synthetic source only.
\subsection{Does simply adding more shrinkage intervals close the gap?}
\label{sec:p4-refinement}
To test whether the observed benefit of the
whole-numerator Bernstein formulation is merely
an artefact of using insufficiently many intervals,
we reran the same fifteen deterministic synthetic
source spectra with $12$, $24$ and $48$
equal-length subintervals. The original
$12$-subinterval run already used source
pair features held fixed across comparisons.
The measured mean improvement of Bernstein
$48$ segments over its own $12$-segment
coverage was just \emph{$+0.0412$ percentage
points} over these fifteen cases, with a
strict increase in only $2/15$.
No case gave a decrease beyond the
stated numerical tolerance in this test.

The result supports a narrow, mechanistic
reading of the earlier synthetic comparisons:
the much larger differences between V23
termwise/Taylor and V25 whole-numerator
Bernstein at identical $12$-segment budgets
cannot be attributed exclusively to
insufficient segmentation of the latter
bound. Exploiting polynomial cancellations
is a different enclosure choice from
simply refining the grid. This does not
prove optimality of Bernstein bounds,
because neither enclosure includes exact
root isolation and each is conditional
on floating-point eigensystems.
The $12$-, $24$- and $48$-segment runs also
do not establish better runtime efficiency;
denser subdivision has a computational cost.
All cases are \emph{synthetic}, not
additional real EEG recordings.


\section{Limitations and independent-study requirement}
These percentages measure \emph{certified
SOURCE pair stability}, NOT new EEG
classification accuracy, AUC improvement
or the fraction of clinically
beneficial decisions. The experiment
uses a single previously explored
three-channel acquisition with nine
source validation groups. Two
different source certificate implementations
on the same fitted model are not
independent held-target experiments.
The Bernstein bounding method has
a floating-point safety margin rather
than a formal outward-directed rounding
guarantee. Earlier numerical problems,
large tangent dimensions and finite
root conditioning remain test cases.
The original unshifted Horner comparison found
additional Bernstein certificates, but the new
midpoint-centered Horner baseline achieved the
same strict rational coverage with lower measured
time than full Bernstein on these 960 intervals.
A corrected 4,320-case exploratory rational sweep found
twenty additional Bernstein signs across several configurations.
A separate exact incremental event aggregator improves timing
on constructed crossing-rich cases, with no stable no-event gain;
no real-input or root-isolation advantage was established.
A publishable method note would
need a substantial independent benchmark
on multiple spectral dimensions and
an explicit complexity/precision tradeoff.
The resulting certificate could be useful
without a new classifier, but its
downstream value must be shown rather
than assumed.

\section{Conclusion}
The six-feature source score path admits
a finite polynomial rank-cell analysis.
Whole-numerator Bernstein bounds
raised the certified source pair coverage
only slightly on the actual Weibo2014
retrospective experiment. The finding
is technical and bounded, and should
not be published as a surrogate
for transfer-performance gains.

\section*{Statements and Declarations}
\textbf{Data/code:} Real source-only
workflow and SHA-256 artifact are
available in the research repository;
independent release and source code
stabilization are not yet complete.
\textbf{Authorship, conflicts,
funding, ethics, AI use:} must be
approved by actual authors.
\bibliographystyle{unsrt}
\bibliography{references_v10}
\end{document}
